\subsection*{Appendix A.1: Proof of Theorem 1 (Optimality of Variance Equalization)}
\label{app:proof_thm1}

\begin{theorem}[Optimality of Variance Equalization]
Under the \textit{convex block error model} (i.e., the MSE $\phi_b(\sigma^2)$ of block $b$ is a convex function of $\sigma^2$), the total MSE satisfies:
\begin{equation}
\min_{\mathbf{R} \in \mathcal{O}(B)} \sum_{b=1}^B \phi_b \left( [\mathbf{R} \boldsymbol{\Sigma}_k \mathbf{R}^\top]_{bb} \right) = \sum_{b=1}^B \phi_b \left( \frac{\mathrm{tr}(\boldsymbol{\Sigma}_k)}{B} \right)
\end{equation}
Equality holds if and only if $\mathrm{diag}(\mathbf{R} \boldsymbol{\Sigma}_k \mathbf{R}^{\top}) = \frac{\mathrm{tr}(\boldsymbol{\Sigma}_k)}{B} \mathbf{1}$.
\end{theorem}

\begin{proof}
\textbf{Step 1: Problem Transformation.}
By the property of orthogonal transformations, the total variance is conserved:
\begin{equation}
\sum_{b=1}^B [\mathbf{R} \boldsymbol{\Sigma}_k \mathbf{R}^\top]_{bb} = \mathrm{tr}(\boldsymbol{\Sigma}_k) = \text{constant}
\end{equation}
Let $\mathbf{d} = \mathrm{diag}(\mathbf{R} \boldsymbol{\Sigma}_k \mathbf{R}^\top)$. The optimization problem can be reformulated as:
\begin{equation}
\min_{\mathbf{d} \in \mathcal{D}} \sum_{b=1}^B \phi_b(d_b) \quad \text{s.t.} \quad \sum_{b=1}^B d_b = \mathrm{tr}(\boldsymbol{\Sigma}_k)
\end{equation}
where $\mathcal{D}$ is the set of achievable diagonal vectors.

\textbf{Step 2: Application of Schur-Horn Theorem.}
The Schur-Horn theorem~\cite{schur_horn} states that $\mathcal{D}$ is exactly the set of vectors \textit{majorized} by the eigenvalues $\boldsymbol{\lambda} = \mathrm{eig}(\boldsymbol{\Sigma}_k)$. That is:
\begin{equation}
\mathcal{D} = \left\{ \mathbf{d} \in \mathbb{R}^B : \sum_{i=1}^j d_{[i]} \leq \sum_{i=1}^j \lambda_{[i]} \text{ for } j=1,\dots,B-1, \text{ and } \sum_{i=1}^B d_i = \sum_{i=1}^B \lambda_i \right\}
\end{equation}
where $d_{[i]}$ denotes the $i$-th largest component of $\mathbf{d}$.

\textbf{Step 3: Convex Optimization and Majorization Theory.}
According to Karamata's inequality for convex functions~\cite{karamata}: if $\phi$ is convex and $\mathbf{d} \prec \mathbf{e}$ ($\mathbf{d}$ is majorized by $\mathbf{e}$), then $\sum \phi(d_i) \le \sum \phi(e_i)$.
The constant vector $\mathbf{c} = \frac{\mathrm{tr}(\boldsymbol{\Sigma}_k)}{B} \mathbf{1}$ is the \textit{most uniform} vector in $\mathcal{D}$ (i.e., $\mathbf{c} \prec \mathbf{d}$ for any $\mathbf{d} \in \mathcal{D}$).
Since $\phi_b$ is convex, the inequality $\sum \phi_b(c) \le \sum \phi_b(d_b)$ holds for all $\mathbf{d} \in \mathcal{D}$.

\textbf{Step 4: Equality Condition.}
Equality holds if and only if:
1) $\mathbf{d} = \mathbf{c}$ (from the equality condition of Karamata's inequality).
2) $\mathbf{c} \in \mathcal{D}$ (from Schur-Horn, since $\mathbf{c}$ is majorized by $\boldsymbol{\lambda}$).
The condition that $\mathbf{c}$ is majorized by $\boldsymbol{\lambda}$ is equivalent to $\lambda_{\min} \le c \le \lambda_{\max}$, which is naturally satisfied as $\frac{\mathrm{tr}(\boldsymbol{\Sigma}_k)}{B}$ lies between the minimum and maximum eigenvalues.

\textbf{Step 5: Constructive Existence.}
By the Schur-Horn theorem, there exists an orthogonal matrix $\mathbf{R}$ such that $\mathrm{diag}(\mathbf{R} \boldsymbol{\Sigma}_k \mathbf{R}^\top) = \mathbf{c}$. This can be constructed via eigenvalue decomposition followed by a series of Givens rotations (as detailed in Algorithm 1) or via a closed-form Householder solution.
Specifically:
1) Perform eigen-decomposition $\boldsymbol{\Sigma}_k = \mathbf{U} \boldsymbol{\Lambda} \mathbf{U}^\top$.
2) Solve for $\mathbf{Q}$ such that $\mathrm{diag}(\mathbf{Q} \boldsymbol{\Lambda} \mathbf{Q}^\top) = \mathbf{c}$.
3) Let $\mathbf{R} = \mathbf{Q}\mathbf{U}^\top$.
\end{proof}

\textbf{Key Insight:} This theorem indicates that under reasonable convex error assumptions (as proven for MXFP4 in Appendix A.2), a rotation configuration that perfectly equalizes variances across all blocks at a given position (i.e., $d_b = \mathrm{tr}(\boldsymbol{\Sigma}_k)/B$ for all $b$) is the global optimal solution. This theoretically underpins the design goal of inter-block rotation: to eliminate the dominance of a few high-variance blocks and distribute the error uniformly.